\nocite{*}
\paragraph{Scope and how to read this chapter.}
This chapter covers all \textbf{49 records} in the supplied 1964--65
catalogue, with \textbf{248 mathematical occurrences}.  All 247 supplied
identifiers are retained; the additional support-of-measure observation
is numbered 1964-6.9.  Dates range from
21 September 1964 to 10 May 1965; the February planning entry is included
without inventing mathematical content.  The sequence follows the supplied
MCCME record order, which is not chronological.  Individual page citations
locate the actual, sometimes interleaved, notebook text.  P000090 starts the
following academic year's Kohn material and is excluded from the present
occurrence count.  Russian transcription is the textual authority; scans
decide consequential doubtful readings.  In each entry, \emph{The notes},
\emph{Statement}, and \emph{Reconstruction} distinguish the record, a precise
modern formulation, and the editor's additions.

\paragraph{The year's mathematical lines.}
The seminar alternates sustained exposition with short announcements and
questions.  Blowups, representation theory, Hilbert kernels, distributions,
Fredholm operators, algebraic jets, and sites supply common language.
Statistical mechanics leads from contour arguments to a walk question;
random matrices lead from circular ensembles to gap probabilities and
spectral traces.  Arnold's geometric optics connects caustics to quantization,
while Novikov's exposition places vector fields on spheres in real
\(K\)-theory.  The final Halmos session collects several decades of
invariant-subspace methods and ends with unusually concrete weighted
operators.  These connections are mathematical and textual; the chapter does
not infer influence merely from two related topics appearing together.

\paragraph{Publication and identification.}
Some announcements have close later matches; others were already published
or circulating.  Dobrushin's phase-separation work requires its correction
as well as the original theorem.  The different-rank Weyl-algebra question
(1964-38.1) has an affirmative distinction through dimension growth; it is
not the fixed-rank endomorphism conjecture.  Sz.-Nagy--Foia\c{s}'s contraction
theorem is available before the Halmos talk, and its linked erratum leaves
that theorem intact.  Domar's later weighted-shift construction resolves the
question of an invariant-subspace chain indexed by all integers, after
adjoining zero and the whole space.  Submission dates are reported as submission dates:
they are not silently identified with first public circulation.  A modern
proof of a recorded idea, such as the Slater-determinant route to Sturm's
theorem, is distinguished from a new theorem discovered in the notebook.

\begin{center}
\begin{tabular}{@{}lrr@{}}\toprule
Class&This chapter&Supplied first run\\\midrule
A&24&22\\
B&22&21\\
C&5&7\\
D&0&0\\
E&160&163\\
F&37&34\\
\midrule Total&248&247\\\bottomrule
\end{tabular}
\end{center}
Here A denotes a close subsequent publication by the speaker or Gelfand;
B a later development by others; C a partly settled precise line; D a
searched but undeveloped line; E prior mathematics; and F an identification
the record cannot sustain.  The separate correctness and match fields are
essential: a familiar theorem may be false in the notebook's literal
wording.  E is an expository classification, not a separately audited
priority claim for every assertion: some record-6 entries lack a verified
pre-talk theorem-level publication date, as their fates expressly state.
No absence of a search hit is used as evidence of an orphan idea.

\paragraph{Questions and bounded development.}
The 17 problem entries include both settled and unresolved questions,
with explicit hypotheses.  The major residual operator questions concern
the Hilbert invariant-subspace problem and Bishop operators.  Weight
classification is reported only to the scope of the verified criteria.
The walk question needs its sampling law specified before a
modern asymptotic question is assigned to it.  Sinai's observation on
spectral multiplicity is not converted into a conjecture absent from the
source.  Recent 2026 claims are reported with their actual verification
level, not treated as accepted resolutions.

One development is provided, for Bishop cyclicity.  It proves rational
cyclicity and a finite-orbit continuity bound, supplies exact approximation
certificates checked by two independent implementations, and derives a
terminating sufficient continued-fraction criterion.  These are bounded
certification and exposition results, not a solution for every irrational
rotation and not a claim of historical novelty; compare
\citet{chamizo2020bishop} and \citet{behani2023bishop}.

\paragraph{Continuity with the supplied index.}
Despite its filename, \texttt{previous\_index.md} contains this same year's
247 first-run items, no earlier-year occurrences, and no existing problem
identifiers.  It is therefore a comparison baseline, not evidence for
cross-year links.  All retained identifiers remain stable; 14 class
assignments change after renewed source checking.  Within-year links join
the two invariant-subspace announcements, repeated gap formulas, and the
successive algebraic-jet discussions.  A genuine earlier-year index can be
joined later without fabricating a history that was not supplied.
